\documentclass[handout,x11names]{beamer}
\usepackage{etex,beamerfoils}
\usepackage{array,pifont,amsmath,latexsym,epsfig,amsfonts,yfonts}
\usepackage{amstext,amsopn,amsxtra,amsgen,amsbsy,amscd,amssymb,dcolumn,graphicx,setspace,pgfplots,booktabs,natbib}
\usepackage[normalem]{ulem}
\usepackage[all]{xy}
\beamertemplatetheoremsnumbered
\setbeamertemplate{caption}[numbered]
\usetheme{Copenhagen}
\definecolor{NewBlue}{RGB}{4, 114, 155}
\definecolor{magenta}{cmyk}{0,1,0,0}
\usecolortheme[named=NewBlue]{structure}
%\usecolortheme[named=SkyBlue4]{structure}


\newcommand{\D}{\displaystyle}
\newcommand{\E}{\begin{eqnarray*}}
\newcommand{\F}{\end{eqnarray*}}
\newcommand{\EE}{\begin{eqnarray}}
\newcommand{\FF}{\end{eqnarray}}
\newcommand{\IZ}{\begin{itemize}}
\newcommand{\ZI}{\end{itemize}}
\newcommand{\EN}{\begin{enumerate}}
\newcommand{\NE}{\end{enumerate}}
\newcommand{\itemc}{\item[$\circ$]}
\newcommand{\IZdash}{\begin{itemize} \renewcommand{\labelitemi}{-}}
\newcommand{\tn}{\textnormal}

% New commands to assist Beamer
\newcommand{\itemi}{\item<1->}
\newcommand{\itemii}{\item<2->}
\newcommand{\itemiii}{\item<3->}
\newcommand{\itemiv}{\item<4->}
\newcommand{\itemv}{\item<5->}
\newcommand{\itemvi}{\item<6->}
\newcommand{\itemvii}{\item<7->}
\newcommand{\itemviii}{\item<8->}
\newcommand{\itemix}{\item<9->}
\newcommand{\itemx}{\item<10->}

\newcommand{\itemci}{\item[$\circ$]<1->}
\newcommand{\itemcii}{\item[$\circ$]<2->}
\newcommand{\itemciii}{\item[$\circ$]<3->}
\newcommand{\itemciv}{\item[$\circ$]<4->}
\newcommand{\itemcv}{\item[$\circ$]<5->}
\newcommand{\itemcvi}{\item[$\circ$]<6->}
\newcommand{\itemcvii}{\item[$\circ$]<7->}
\newcommand{\itemcviii}{\item[$\circ$]<8->}
\newcommand{\itemcix}{\item[$\circ$]<9->}
\newcommand{\itemcx}{\item[$\circ$]<10->}


\newcommand{\fr}{\frame}
\newcommand{\ft}{\frametitle}
\newcommand{\ons}{\onslide}
\newcommand{\tc}{\textcolor}

\pgfplotsset{height =7cm,compat=1.7,width=7cm,
    standard/.style={
    	unit vector ratio*=1 1 1,
        axis x line=middle,
        axis y line=middle,
        enlarge x limits=0.15,
        enlarge y limits=0.15,
        every axis x label/.style={at={(current axis.right of origin)},anchor=north west},
        every axis y label/.style={at={(current axis.above origin)},anchor=north east}
    }
}
\newcommand{\drawge}{-- (rel axis cs:1,0) -- (rel axis cs:1,1) -- (rel axis cs:0,1) \closedcycle}
\newcommand{\drawle}{-- (rel axis cs:1,1) -- (rel axis cs:1,0) -- (rel axis cs:0,0) \closedcycle}

\usetikzlibrary{patterns}
\usetikzlibrary{arrows}
\usetikzlibrary{shapes}

\renewcommand{\bibsection}{}

% Watermark example: https://texblog.org/2012/02/17/watermarks-draft-review-approved-confidential/
\usepackage{draftwatermark}
\SetWatermarkText{This content is protected and may not be shared uploaded or distributed.}
\SetWatermarkScale{0.25}
\SetWatermarkAngle{0}
\SetWatermarkColor[gray]{1.0}
\setbeamercolor{background canvas}{bg=}%transparent canvas

\title[Deterministic Time Series Models\hspace{2em}\insertframenumber/
\inserttotalframenumber]{Deterministic Time Series Models}
\author[Brian C. Jenkins - Computational Macroeconomics ]{Brian C.~Jenkins \\ ~ \\Econ 126: Computational Macroeconomics\\ ~ \\University of California, Irvine}

\begin{document}

\frame{\maketitle}



\fr{
\ft{Introduction}

\IZ
\item Macroeconomic data evolve over time and there is good evidence that:
	\EN
	\itemii Historical events matter for the current economy
	\itemiii Current events will matter for the future economy
	\NE
	
\

\itemiv We model the \emph{dynamic} nature of the economy using \emph{time series} models

\

\itemv A  \textbf{time series} model specifies how a variable or collection of variables are determined as a function of time.
\ZI
}

\fr{
\ft{Introduction}

\IZ
\item A  \textbf{time series} model specifies how a variable or collection of variables are determined as a function of time.

\

\itemii Time series models can be:

	\EN
	\itemiii \textbf{Deterministic}: non-random and purely mechanical
	\itemiv \textbf{Stochastic}: Well-defined mathematical structure, but with random elements
	\NE
	
\

\itemv Deterministic models are well-suited for modeling long-run aspects of the economy

\

\itemvi Stochastic models are ideal for modeling business cycle fluctuations because the consensus is that cycles are caused by unpredictable disturbances
\ZI
}


\fr{
\ft{Discrete Versus Continuous Time}

\IZ
\item Let $y$ denote a variable that takes on the value $y_t$ at date $t$.

\


\itemii If $t$ takes on values from a countable sequence, e.g., $t\in [0,1,2,\ldots)$, then $y_t$ is a   \textbf{discrete time} variable.

\

\itemiii Otherwise, if $t$ takes on values from an uncountable sequence, e.g., $t\in [0,\infty)$, then $y_t$ is a \textbf{continuous time} variable.

\

\itemiv We will focus exclusively on discrete time models.
\ZI

}



\fr{
\ft{First-Order Difference Equations}

\IZ
\item Suppose that the variable $y_t$ is determined by a linear function of $y_{t-1}$ and some other exogenously given variable $w_t$:
\begin{align}
y_{t} & = \rho y_{t-1} + w_t,  \label{eqn:first_diff}
\end{align}
where $\rho$ is some constant.

\

\itemii Equation (\ref{eqn:first_diff}) is an example of a \textbf{linear first-order difference equation}.


\ZI

}

\fr{
\ft{Examples}
\begin{exampleblock}{Example: Compounding Interest}

\IZ
\item Suppose that you have an initial balance of $b_0$ dollars in an account that pays an interest rate $i$ per compounding period.

\itemii Your period $t$ balance depends on your period $t-1$ balance:
\begin{align}
b_{t} & =  \left(1+i\right) b_{t-1}. \label{eqn:compound_interst}
\end{align}

\

\itemiii Equation (\ref{eqn:compound_interst}) is a linear first-order difference equation in the same form as Equation (\ref{eqn:first_diff}). You can see this by setting $y_t = b_t$, $\rho=1+i$, and $w_t=0$ in Equation (\ref{eqn:first_diff}).
\ZI
\end{exampleblock}

}


\fr{
\ft{Examples}
\begin{exampleblock}{Example: Physical Capital Accumulation}

\IZ
\itemi In the Solow growth model, the law of motion for physical capital is:
\begin{align}
K_{t+1} & =  I_t + (1-\delta)K_t, \label{eqn:capital_accumulation}
\end{align}

where $K_t$ is the capital stock in period $t$, $\delta$ is the rate of capital depreciation, and $I_t$ is investment in new capital

\

\itemii Treating investment $I_t$ as exogenous, Equation (\ref{eqn:capital_accumulation}) is a linear first-order difference equation in the same form as Equation (\ref{eqn:first_diff}). You can see this by setting $y_t = K_{t+1}$, $\rho=1-\delta$, and $w_t=I_t$ in Equation (\ref{eqn:first_diff}).
\ZI
\end{exampleblock}

}



\end{document}